\section{Related work}
\label{sec:related_work}

\subsection{Physics-informed learning with scarce observations}

PINNs combine observations with differential-equation, initial-condition, and boundary-condition residuals, so the governing model supplies supervision at points for which no reference solution is available \cite{Raissi2019,Karniadakis2021}. This distinction has motivated their use when measurements or simulator evaluations are scarce. Sun et al. constructed parametric fluid-flow surrogates using equation-based training without simulation labels \cite{Sun2020}, while He et al. jointly reconstructed conductivity, hydraulic head, and concentration from sparse measurements in a subsurface data-assimilation problem \cite{He2020}. In the latter study, the physics-informed networks were more accurate than otherwise comparable data-driven neural networks as the measurement set became sparse. Physics-informed learning has also been coupled with sparse regression to identify governing equations from scarce and noisy observations \cite{Chen2021}. These results support the premise that known physics can reduce dependence on labelled states, but they address data-free solution, inverse reconstruction, or comparison with a data-only neural network rather than the label-budget ranking against conventional non-neural surrogates studied here.

The distinction between labels and residual points is important in interpreting data efficiency. A collocation point does not require a measurement or a reference-solver state, but evaluating the differential operator there is not computationally free. Accordingly, we use ``small data'' exclusively for the number of labelled solution states, report the residual budget separately, and measure fitting time. This avoids describing a PINN as computationally data-free merely because its additional supervision is generated from a known PDE.

\subsection{Comparisons with numerical and statistical alternatives}

The strongest alternative depends on the task. For direct solution of a known PDE, a mature discretization can be considerably more efficient than training a PINN. In a systematic comparison across several canonical PDEs, Grossmann et al. found that finite elements were faster at a given accuracy, although a trained PINN could sometimes be evaluated rapidly \cite{Grossmann2024}. Optimization studies likewise show that PINNs can fail even when their architecture is expressive: differential operators can make the training landscape ill-conditioned, and curriculum or time-decomposed strategies may be necessary \cite{Krishnapriyan2021}. The PINNacle benchmark subsequently compared approximately ten PINN methods on more than twenty PDEs and confirmed that performance depends strongly on the equation, geometry, and training strategy \cite{Hao2024}. These findings preclude a universal claim that PINNs replace established numerical solvers.

Our comparison concerns a different decision: after a high-fidelity solver or experiment has supplied only (N) scattered states, which surrogate reconstructs the remaining space--time field most accurately? XGB, radial-basis SVR, and GP are relevant because they provide nonlinear regression with very different inductive biases and are effective in low-dimensional, small-sample settings \cite{chen2016xgboost,drucker1997support,rasmussen2006gaussian}. Gaussian processes have also been used directly for function approximation and PDE solution; Pang et al. found that GP-type models were particularly competitive for smooth prototype problems \cite{Pang2019}. In the present study, however, GP is intentionally a data-only statistical baseline rather than a physics-informed GP. The numerical solver defines the reference field and the cost of obtaining labels; it is not treated as a label-matched surrogate. This separation lets us test whether residual supervision changes the learning curve, without conflating surrogate accuracy with the cost and accuracy of the reference discretization.

Existing PINN benchmarks mainly compare PINN variants with one another or compare a neural PDE solver with a discretization \cite{Grossmann2024,Hao2024}. Conversely, conventional-surrogate studies often do not give one model access to the governing residual. A controlled comparison of these two information sets over nested label budgets, paired random designs, and separate interpolation and future-time regions is therefore less well characterized. Our goal is not to assert that the comparison is information-equivalent---the additional known physics is precisely the PINN's proposed advantage---but to hold the expensive information, namely labelled states, fixed.

\subsection{PINNs for porous-media and two-phase flow}

Porous-media flow is a natural small-data application because heterogeneous coefficients, coupled state variables, and moving saturation fronts make dense simulation or measurement campaigns expensive. Physics-informed subsurface studies have reconstructed multiple coupled fields from sparse observations \cite{He2020}. For two-phase flow, Hanna et al. showed that residual-based adaptive sampling is important for resolving a moving front and outperforms fixed collocation and standard residual-based refinement in their examples \cite{Hanna2022}. Yan et al. developed an enriched PINN for heterogeneous and fractured two-phase systems, using finite-volume spatial residuals, hard initial and boundary constraints, and a specialized architecture to handle discontinuous material properties \cite{Yan2024}.

Recent work has also examined pressure and saturation in a unified network. Kenzhebek et al. reported that a standard coupled PINN struggled with the hyperbolic saturation shock and that artificial diffusion, combined with optimized loss weighting, was needed for reliable front capture in their formulation \cite{Kenzhebek2025}. This observation is directly relevant to our choice of water saturation as the primary flow target: pressure is comparatively smooth, whereas saturation exposes whether a surrogate preserves a sharp advancing displacement structure. It also motivates reporting pressure and saturation separately rather than allowing the easier pressure field to dominate an aggregate score.

The present benchmarks complement those studies in two ways. First, the reference states are generated by a corrected conservative finite-volume scheme, and exactly the same nested labels are given to every surrogate. Second, the central independent variable is the labelled budget rather than a particular front-capturing modification. We retain a common PINN backbone across the tests and ask where its equation-based regularization changes the ranking relative to XGB, SVR, and GP. We therefore do not claim to supersede specialized two-phase PINN solvers; instead, we study the practical surrogate-selection question under controlled scarcity.

\subsection{Balancing composite PINN objectives}

The data, initial/boundary, and equation terms of a PINN form a multi-objective optimization problem. Their numerical scales and gradient magnitudes can differ substantially. Wang et al. connected this imbalance to gradient-flow pathologies and proposed gradient-statistics-based annealing \cite{Wang2022}. Other approaches adapt the objective in different ways. Self-adaptive PINNs learn pointwise weights that emphasize difficult locations \cite{McClenny2023}; inverse-Dirichlet weighting uses gradient statistics to reduce scale imbalance and catastrophic interference \cite{Maddu2022}; and ReLoBRaLo balances objectives through their relative progress with randomized historical lookback \cite{Bischof2025}. Curriculum formulations instead alter when or where the physical constraints are presented \cite{Krishnapriyan2021}. More generally, projected-gradient methods such as PCGrad modify mutually conflicting task gradients before combining them \cite{Yu2020}.

H-PINN and HO-PINN occupy a distinct, narrower position in this design space. H-PINN does not attempt to equalize all loss magnitudes or learn one weight per collocation point. It encodes an explicit priority order---labelled data, then initial/boundary conditions, then the PDE residual---and uses exponential gates to activate a lower-priority objective as the preceding objectives approach specified target scales. HO-PINN retains this hierarchy and multiplies each lower-priority gate by the fraction of its gradient lying outside the span of preceding gradients. This is related to gradient-interference methods in its use of gradient geometry, but it is not PCGrad: HO-PINN does not replace each task gradient by a projected gradient or remove only negative pairwise conflicts. It uses a detached scalar novelty factor, with a positive floor, to modulate the hierarchical objective before the ordinary backward pass. The exact formulas and screened target scales are given in Section~\ref{sec:method} and \ref{app:tuning}.

\subsection{Position of the present study}

Prior work establishes separately that physics can guide learning from sparse observations, that conventional and numerical methods remain strong competitors, that two-phase saturation fronts are difficult for PINNs, and that composite-loss optimization materially affects accuracy. The empirical gap addressed here is their intersection: a label-matched, training-only-tuned comparison against strong non-neural surrogates over a common sequence of scarce-data budgets, with pressure and water saturation assessed separately and interpolation distinguished from physics-assisted temporal forecasting. The resulting claim is deliberately conditional. We seek the budgets and physical regimes in which PINNs are advantageous, the regimes in which a conventional surrogate catches up, and the computational price paid for the additional equation-based supervision.
